\documentclass[10pt,a4paper]{amsart}
\usepackage{amsmath,amssymb,amsthm,graphicx}
\usepackage{verbatim}
\newenvironment{reference}{\comment}{\endcomment}
%\newenvironment{reference}{}{}
\newenvironment{slogan}{\comment}{\endcomment}
\newenvironment{history}{\comment}{\endcomment}

% For commutative diagrams we use Xy-pic
\usepackage[all]{xy}

% We use 2cell for 2-commutative diagrams.
\xyoption{2cell}
% Theorem environments.
%
\theoremstyle{plain}
\newtheorem{theorem}[subsection]{Theorem}
\newtheorem{proposition}[subsection]{Proposition}
\newtheorem{lemma}[subsection]{Lemma}

\theoremstyle{definition}
\newtheorem{definition}[subsection]{Definition}
\newtheorem{example}[subsection]{Example}
\newtheorem{exercise}[subsection]{Exercise}
\newtheorem{situation}[subsection]{Situation}

\theoremstyle{remark}
\newtheorem{remark}[subsection]{Remark}
\newtheorem{remarks}[subsection]{Remarks}

\numberwithin{equation}{subsection}

% Macros
%
\def\lim{\mathop{\mathrm{lim}}\nolimits}
\def\colim{\mathop{\mathrm{colim}}\nolimits}
\def\Spec{\mathop{\mathrm{Spec}}}
\def\Hom{\mathop{\mathrm{Hom}}\nolimits}
\def\Ext{\mathop{\mathrm{Ext}}\nolimits}
\def\SheafHom{\mathop{\mathcal{H}\!\mathit{om}}\nolimits}
\def\SheafExt{\mathop{\mathcal{E}\!\mathit{xt}}\nolimits}
\def\Sch{\mathit{Sch}}
\def\Mor{\mathop{\mathrm{Mor}}\nolimits}
\def\Ob{\mathop{\mathrm{Ob}}\nolimits}
\def\Sh{\mathop{\mathit{Sh}}\nolimits}
\def\NL{\mathop{N\!L}\nolimits}
\def\CH{\mathop{\mathrm{CH}}\nolimits}
\def\proetale{{pro\text{-}\acute{e}tale}}
\def\etale{{\acute{e}tale}}
\def\QCoh{\mathit{QCoh}}
\def\Ker{\mathop{\mathrm{Ker}}}
\def\Im{\mathop{\mathrm{Im}}}
\def\Coker{\mathop{\mathrm{Coker}}}
\def\Coim{\mathop{\mathrm{Coim}}}

% Boxtimes
%
\DeclareMathSymbol{\boxtimes}{\mathbin}{AMSa}{"02}

%
% Macros for moduli stacks/spaces
%
\def\QCohstack{\mathcal{QC}\!\mathit{oh}}
\def\Cohstack{\mathcal{C}\!\mathit{oh}}
\def\Spacesstack{\mathcal{S}\!\mathit{paces}}
\def\Quotfunctor{\mathrm{Quot}}
\def\Hilbfunctor{\mathrm{Hilb}}
\def\Curvesstack{\mathcal{C}\!\mathit{urves}}
\def\Polarizedstack{\mathcal{P}\!\mathit{olarized}}
\def\Complexesstack{\mathcal{C}\!\mathit{omplexes}}
% \Pic is the operator that assigns to X its picard group, usage \Pic(X)
% \Picardstack_{X/B} denotes the Picard stack of X over B
% \Picardfunctor_{X/B} denotes the Picard functor of X over B
\def\Pic{\mathop{\mathrm{Pic}}\nolimits}
\def\Picardstack{\mathcal{P}\!\mathit{ic}}
\def\Picardfunctor{\mathrm{Pic}}
\def\Deformationcategory{\mathcal{D}\!\mathit{ef}}
\usepackage{cleveref}
\crefname{originalsection}{original section}{original sections}
\Crefname{originalsection}{Original section}{Original sections}
\makeatletter
\newlabel{schemes-section-quasi-compact}{{01K2}{source}{}{sourcepointer-1}{}}
\newlabel{schemes-section-quasi-compact@cref}{{[originalsection][0][]01K2}{[1][0][]source}{}{}{}}
\makeatother
\crefname{originalsection}{original section}{original sections}
\Crefname{originalsection}{Original section}{Original sections}
\makeatletter
\newlabel{section-representable}{{03HA}{source}{}{sourcepointer-2}{}}
\newlabel{section-representable@cref}{{[originalsection][0][]03HA}{[1][0][]source}{}{}{}}
\makeatother
\crefname{originaldefinition}{original definition}{original definitions}
\Crefname{originaldefinition}{Original definition}{Original definitions}
\makeatletter
\newlabel{spaces-properties-definition-quasi-compact}{{03E3}{source}{}{sourcepointer-3}{}}
\newlabel{spaces-properties-definition-quasi-compact@cref}{{[originaldefinition][0][]03E3}{[1][0][]source}{}{}{}}
\makeatother
\title{Quasi-compactness of a representable morphism of algebraic spaces tested after base change from quasi-compact spaces}
\author{The Stacks Project Authors}
\date{stacks-project:spaces-morphisms.tex:local-snapshot-20261012; GNU FDL 1.2 or later}
\begin{document}
\maketitle
\setcounter{subsection}{1}
\noindent\textit{Editorial note.} All delivered mathematical text is the author's own, copied from the pinned Stacks Project chapter file spaces-morphisms.tex; no statement or proof step was rewritten, reordered or removed.\par\medskip
\noindent\textit{Editorial note.} The proof is short but complete in both directions; the second direction uses the model case of an affine open of a scheme, which is exactly what the definition of quasi-compactness of a morphism of algebraic spaces asks for.\par\medskip
\noindent\textit{Editorial note.} The three references of the counted spans are recorded as non-load-bearing navigation pointers; the argument that consumes them stays inside the counted proof span.\par\medskip
\noindent\textit{Editorial note.} Preamble: the three native excerpt spans of sources/stacks-spaces-morphisms/preamble.tex (lines 8-13, 14-18 and 154-217). The chapter file itself carries no preamble text.\par\medskip
\noindent\textit{Allowed standard prerequisites (editorial).}\begin{itemize}
\item Quasi-compactness of an algebraic space tested by a surjective etale cover. Properties of Spaces, Definition spaces-properties-definition-quasi-compact, used at proof lines 1226-1230 and 1234-1238. Reduces quasi-compactness of the base change to a scheme cover.
\item Quasi-compactness in the representable sense. Section section-representable, used at proof lines 1218-1219 and 1231-1237. Fixes the meaning of condition (1) and gives quasi-compactness of U x\_Y X -> U.
\item Quasi-compactness of a morphism of schemes tested on affine opens of the target. Schemes, Section schemes-section-quasi-compact, used at proof lines 1241-1246. Converts quasi-compactness of the preimages into quasi-compactness of p.
\end{itemize}\medskip
\ifdefined\hypertarget\hypertarget{sourcepointer-1}{}\fi
\noindent\textit{Source reference.} Original section 01K2: focused-seeds/AGP/sources/stacks-snapshot-20261011/tags.txt, source lines 2035--2035. Navigation pointer into the Stacks Project; the standard fact it supplies is applied inside the counted proof span and the argument resting on it stays there.\par\medskip
\ifdefined\hypertarget\hypertarget{sourcepointer-2}{}\fi
\noindent\textit{Source reference.} Original section 03HA: focused-seeds/AGP/sources/stacks-snapshot-20261011/tags.txt, source lines 4559--4559. Navigation pointer into the Stacks Project; the standard fact it supplies is applied inside the counted proof span and the argument resting on it stays there.\par\medskip
\ifdefined\hypertarget\hypertarget{sourcepointer-3}{}\fi
\noindent\textit{Source reference.} Original definition 03E3: focused-seeds/AGP/sources/stacks-snapshot-20261011/tags.txt, source lines 4442--4442. Navigation pointer into the Stacks Project; the standard fact it supplies is applied inside the counted proof span and the argument resting on it stays there.\par\medskip
\begin{lemma}
\label{lemma-characterize-representable-quasi-compact}
Let $S$ be a scheme.
Let $f : X \to Y$ be a representable morphism of algebraic spaces over $S$.
The following are equivalent:
\begin{enumerate}
\item $f$ is quasi-compact
(in the sense of Section \ref{section-representable}), and
\item for every quasi-compact algebraic space $Z$ and any morphism
$Z \to Y$ the algebraic space $Z \times_Y X$ is quasi-compact.
\end{enumerate}
\end{lemma}

\begin{proof}
Assume (1), and let $Z \to Y$ be a morphism of algebraic spaces with
$Z$ quasi-compact. By
Properties of Spaces,
Definition \ref{spaces-properties-definition-quasi-compact}
there exists a quasi-compact scheme $U$ and a surjective \'etale
morphism $U \to Z$. Since $f$ is representable and quasi-compact
we see by definition that $U \times_Y X$ is a scheme, and that
$U \times_Y X \to U$ is quasi-compact. Hence $U \times_Y X$ is
a quasi-compact scheme. The morphism $U \times_Y X \to Z \times_Y X$
is \'etale and surjective (as the base change of the representable
\'etale and surjective morphism $U \to Z$, see
Section \ref{section-representable}).
Hence by definition $Z \times_Y X$ is quasi-compact.

\medskip\noindent
Assume (2). Let $Z \to Y$ be a morphism, where $Z$ is a scheme.
We have to show that $p : Z \times_Y X \to Z$ is quasi-compact.
Let $U \subset Z$ be affine open. Then $p^{-1}(U) = U \times_Y Z$
and the scheme $U \times_Y Z$ is quasi-compact by assumption (2).
Hence $p$ is quasi-compact, see
Schemes, Section \ref{schemes-section-quasi-compact}.
\end{proof}

\end{document}
